\chapter*{前言}
\BookTopBookmark{前言}{book-preface}
\markboth{前言}{}

\section*{本书目标}

给定一个对象，我们怎样描述它的结构，判断它具有哪些性质，并据此进行计算？
集合、函数、向量与图提供不同的描述方式；极限、测度与概率使这些描述能够处理逼近、
连续变化与随机性；优化、微分方程和控制则进一步研究解怎样求得、状态怎样演化，
以及行动怎样影响结果。本书围绕这些问题整理数学概念、推导与计算方法。

这份笔记从《机器学习笔记》的数学部分独立出来，供需要系统梳理相关数学工具的读者阅读。
各章在具体对象和问题中引入定义，说明结论依赖的条件，并通过证明、算例和反例检查其适用范围。
阅读时可以沿正文顺序建立基础，也可以从当前问题出发查阅相应章节，再补充所需的前置内容。

\section*{内容与阅读}

当前正文保留迁移时的五卷结构。卷一从集合、实数与函数进入微积分、向量矩阵、图与基础概率统计；
卷二研究度量、测度、概率空间与函数空间，进一步说明极限、积分和逼近所依赖的结构；
卷三讨论优化、状态演化、信号表示与估计；卷四连接流形、概率分布、统计推断与因果识别；
卷五研究决策、控制与博弈中的价值、策略和相互作用。各卷独立构建，章节之间按需引用。

数学准备的十卷重组方案正在单独调整，当前尚未据此迁移正文。本文沿用现有卷名、章节和推导，
最新方案与正文提要分别保存在仓库的草案目录和大纲目录中，便于继续核对内容归属与阅读顺序。

正文使用其他数学卷的结论时，会在当前问题中交代所需关系与条件，再给出可查阅的位置。
少量机器学习应用作为延伸阅读，链接指向《机器学习笔记》的固定版本；本书的数学正文与编译过程独立维护。
